% --- Archon physics formalization source begin ---
% archon:physics
% archon:covers PhyXMiniProblems/problem_phyx_mini_0570.lean
% archon:source-report reports/phyx_mini/problem_phyx_mini_0570.source.json

\chapter{Physics problem phyx\_mini\_0570}
\label{ch:PhyXMiniProblems_problem_phyx_mini_0570}

\paragraph{Problem source.}
\#\# Physical scenario

Consider a particle of kinetic energy \$K\$ approaching the step function of the figure from the left, where the potential barrier steps from 0 to \$V\_0\$ at \$x = 0\$.

\#\# Question

Calculate the penetration distance for a 5-eV electron approaching a step barrier of 10 eV.

\#\# Answer choices

- A: A: 0.55nm

- B: B: 0.95nm

- C: C: 0.044nm

- D: D: 10.2nm

\#\# Figure caption (auxiliary; use the image as primary evidence)

The image is a graph depicting a potential step in a quantum mechanics context. It features:

- **Axes**: The horizontal axis is labeled "Position (x)" and the vertical axis is labeled "V(x)" for potential energy.
- **Potential Step**: At position \textbackslash{}( x = 0 \textbackslash{}), there is a vertical line indicating a step in potential energy. The potential is \textbackslash{}( V\_0 \textbackslash{}) for \textbackslash{}( x > 0 \textbackslash{}) and 0 for \textbackslash{}( x < 0 \textbackslash{}).
- **Line Segments**: Horizontal blue lines represent potential energy levels. The left side is at zero potential, and the right side is at potential \textbackslash{}( V\_0 \textbackslash{}).
- **Particle**: An arrow with a dashed line pointing right indicates the movement of a particle approaching the step, labeled "Particle."
- **Energy Level**: A line and label "E" represent the particle's energy, positioned below \textbackslash{}( V\_0 \textbackslash{}).
  
The diagram illustrates a scenario often used to explain quantum tunneling or the behavior of a particle encountering a potential step.

\#\# Dataset metadata

Quantum Phenomena; Physical Model Grounding Reasoning, Numerical Reasoning

\paragraph{Recorded answer/context.}
C: C: 0.044nm

\paragraph{Figure/image path.}
/root/proposal\_for\_physic/science-mango/phyx\_mini\_run/phyx\_data/test\_image/570.png

\paragraph{Formalization target.}
create a compiling Lean file with sorry bodies at `PhyXMiniProblems/problem\_phyx\_mini\_0570.lean`. The Lean declarations must preserve the physical quantities, dimensions or dimensional roles, figure labels, governing-law hypotheses, and final relation expressed by this problem.
Use Mathlib/Physlib names found through LeanExplore where available. If a physics API is missing, introduce faithful local abstractions rather than scalar placeholder aliases.

\begin{theorem}[Physics formalization target]
\label{thm:physics:phyx_mini_0570:target}
\lean{PhyXMiniProblems.ProblemPhyXMini0570.problem_phyx_mini_0570}
\uses{def:physics:phyx-mini-0570:phyxminiproblems-problemphyxmini0570-electronpotentialstepsetup, def:physics:phyx-mini-0570:phyxminiproblems-problemphyxmini0570-matcheselectronstepscenario, def:physics:phyx-mini-0570:phyxminiproblems-problemphyxmini0570-matchesproblemnumericalreadouts, def:physics:phyx-mini-0570:phyxminiproblems-problemphyxmini0570-matchesupwardpotentialstepprofile, def:physics:phyx-mini-0570:phyxminiproblems-problemphyxmini0570-matchessuppliedpotentialstepfigure, def:physics:phyx-mini-0570:phyxminiproblems-problemphyxmini0570-usesstandardelectronreferencedata, def:physics:phyx-mini-0570:phyxminiproblems-problemphyxmini0570-hasphysicalunderstepparameters, def:physics:phyx-mini-0570:phyxminiproblems-problemphyxmini0570-satisfiesevanescentprobabilitydensitylaw, def:physics:phyx-mini-0570:phyxminiproblems-problemphyxmini0570-isprobabilitydensityefoldingdistance, def:physics:phyx-mini-0570:phyxminiproblems-problemphyxmini0570-agreeswhenroundedtothreedecimalnanometers, def:physics:phyx-mini-0570:phyxminiproblems-problemphyxmini0570-isuniquenearestpenetrationdistancechoice}
The assigned autoformalize agent should translate this physics problem into Lean declarations in the covered file, with theorem and lemma proof bodies written as `by sorry`.
\end{theorem}
\begin{proof}
This is an autoformalization task, not a proof task. Produce statements that can later be proved by the physics prover without weakening the physical model.
\end{proof}

% --- Archon declaration topology begin ---
\section{Declaration topology}
\begin{definition}[Length Quantity]
  \label{def:physics:phyx-mini-0570:phyxminiproblems-problemphyxmini0570-lengthquantity}
  \lean{PhyXMiniProblems.ProblemPhyXMini0570.LengthQuantity}
  A nonnegative, unit-independent physical length
\end{definition}
\begin{proof}
  This is the declared model definition of Length Quantity.
\end{proof}
\begin{definition}[Signed Length Quantity]
  \label{def:physics:phyx-mini-0570:phyxminiproblems-problemphyxmini0570-signedlengthquantity}
  \lean{PhyXMiniProblems.ProblemPhyXMini0570.SignedLengthQuantity}
  A signed, unit-independent coordinate along the scattering axis
\end{definition}
\begin{proof}
  This is the declared model definition of Signed Length Quantity.
\end{proof}
\begin{definition}[Mass Quantity]
  \label{def:physics:phyx-mini-0570:phyxminiproblems-problemphyxmini0570-massquantity}
  \lean{PhyXMiniProblems.ProblemPhyXMini0570.MassQuantity}
  A nonnegative, unit-independent physical mass
\end{definition}
\begin{proof}
  This is the declared model definition of Mass Quantity.
\end{proof}
\begin{definition}[Inverse Length Quantity]
  \label{def:physics:phyx-mini-0570:phyxminiproblems-problemphyxmini0570-inverselengthquantity}
  \lean{PhyXMiniProblems.ProblemPhyXMini0570.InverseLengthQuantity}
  A nonnegative inverse length, used for the evanescent decay constant
\end{definition}
\begin{proof}
  This is the declared model definition of Inverse Length Quantity.
\end{proof}
\begin{definition}[action Dimension]
  \label{def:physics:phyx-mini-0570:phyxminiproblems-problemphyxmini0570-actiondimension}
  \lean{PhyXMiniProblems.ProblemPhyXMini0570.actionDimension}
  The physical dimension of action, mass * length\textasciicircum{}2 / time
\end{definition}
\begin{proof}
  This is the declared model definition of action Dimension.
\end{proof}
\begin{definition}[Action Quantity]
  \label{def:physics:phyx-mini-0570:phyxminiproblems-problemphyxmini0570-actionquantity}
  \lean{PhyXMiniProblems.ProblemPhyXMini0570.ActionQuantity}
  \uses{def:physics:phyx-mini-0570:phyxminiproblems-problemphyxmini0570-actiondimension}
  A nonnegative, unit-independent physical action
\end{definition}
\begin{proof}
  This is the declared model definition of Action Quantity.
\end{proof}
\begin{definition}[energy In Joules]
  \label{def:physics:phyx-mini-0570:phyxminiproblems-problemphyxmini0570-energyinjoules}
  \lean{PhyXMiniProblems.ProblemPhyXMini0570.energyInJoules}
  Coherent-SI readout of an energy, in joules
\end{definition}
\begin{proof}
  This is the declared model definition of energy In Joules.
\end{proof}
\begin{definition}[energy In Electron Volts]
  \label{def:physics:phyx-mini-0570:phyxminiproblems-problemphyxmini0570-energyinelectronvolts}
  \lean{PhyXMiniProblems.ProblemPhyXMini0570.energyInElectronVolts}
  \uses{def:physics:phyx-mini-0570:phyxminiproblems-problemphyxmini0570-energyinjoules}
  Read an energy in electron volts using Physlib's calibrated electron volt
\end{definition}
\begin{proof}
  This is the declared model definition of energy In Electron Volts.
\end{proof}
\begin{definition}[mass In Kilograms]
  \label{def:physics:phyx-mini-0570:phyxminiproblems-problemphyxmini0570-massinkilograms}
  \lean{PhyXMiniProblems.ProblemPhyXMini0570.massInKilograms}
  \uses{def:physics:phyx-mini-0570:phyxminiproblems-problemphyxmini0570-massquantity}
  Read a mass in kilograms
\end{definition}
\begin{proof}
  This is the declared model definition of mass In Kilograms.
\end{proof}
\begin{definition}[action In Joule Seconds]
  \label{def:physics:phyx-mini-0570:phyxminiproblems-problemphyxmini0570-actioninjouleseconds}
  \lean{PhyXMiniProblems.ProblemPhyXMini0570.actionInJouleSeconds}
  \uses{def:physics:phyx-mini-0570:phyxminiproblems-problemphyxmini0570-actionquantity}
  Read an action in coherent SI units, joule-seconds
\end{definition}
\begin{proof}
  This is the declared model definition of action In Joule Seconds.
\end{proof}
\begin{definition}[position In Meters]
  \label{def:physics:phyx-mini-0570:phyxminiproblems-problemphyxmini0570-positioninmeters}
  \lean{PhyXMiniProblems.ProblemPhyXMini0570.positionInMeters}
  \uses{def:physics:phyx-mini-0570:phyxminiproblems-problemphyxmini0570-signedlengthquantity}
  Read a signed position in metres
\end{definition}
\begin{proof}
  This is the declared model definition of position In Meters.
\end{proof}
\begin{definition}[length Readout]
  \label{def:physics:phyx-mini-0570:phyxminiproblems-problemphyxmini0570-lengthreadout}
  \lean{PhyXMiniProblems.ProblemPhyXMini0570.lengthReadout}
  \uses{def:physics:phyx-mini-0570:phyxminiproblems-problemphyxmini0570-lengthquantity}
  Read a physical length in a selected length unit
\end{definition}
\begin{proof}
  This is the declared model definition of length Readout.
\end{proof}
\begin{definition}[length In Meters]
  \label{def:physics:phyx-mini-0570:phyxminiproblems-problemphyxmini0570-lengthinmeters}
  \lean{PhyXMiniProblems.ProblemPhyXMini0570.lengthInMeters}
  \uses{def:physics:phyx-mini-0570:phyxminiproblems-problemphyxmini0570-lengthquantity, def:physics:phyx-mini-0570:phyxminiproblems-problemphyxmini0570-lengthreadout}
  Read a physical length in metres
\end{definition}
\begin{proof}
  This is the declared model definition of length In Meters.
\end{proof}
\begin{definition}[length In Nanometers]
  \label{def:physics:phyx-mini-0570:phyxminiproblems-problemphyxmini0570-lengthinnanometers}
  \lean{PhyXMiniProblems.ProblemPhyXMini0570.lengthInNanometers}
  \uses{def:physics:phyx-mini-0570:phyxminiproblems-problemphyxmini0570-lengthquantity, def:physics:phyx-mini-0570:phyxminiproblems-problemphyxmini0570-lengthreadout}
  Read a physical length in nanometres
\end{definition}
\begin{proof}
  This is the declared model definition of length In Nanometers.
\end{proof}
\begin{definition}[inverse Length In Inverse Meters]
  \label{def:physics:phyx-mini-0570:phyxminiproblems-problemphyxmini0570-inverselengthininversemeters}
  \lean{PhyXMiniProblems.ProblemPhyXMini0570.inverseLengthInInverseMeters}
  \uses{def:physics:phyx-mini-0570:phyxminiproblems-problemphyxmini0570-inverselengthquantity}
  Read an inverse-length quantity in inverse metres
\end{definition}
\begin{proof}
  This is the declared model definition of inverse Length In Inverse Meters.
\end{proof}
\begin{definition}[Particle Species]
  \label{def:physics:phyx-mini-0570:phyxminiproblems-problemphyxmini0570-particlespecies}
  \lean{PhyXMiniProblems.ProblemPhyXMini0570.ParticleSpecies}
  Particle species distinguished by the problem statement
\end{definition}
\begin{proof}
  This is the declared model definition of Particle Species.
\end{proof}
\begin{definition}[Potential Region]
  \label{def:physics:phyx-mini-0570:phyxminiproblems-problemphyxmini0570-potentialregion}
  \lean{PhyXMiniProblems.ProblemPhyXMini0570.PotentialRegion}
  The two constant-potential regions separated by x = 0
\end{definition}
\begin{proof}
  This is the declared model definition of Potential Region.
\end{proof}
\begin{definition}[Horizontal Direction]
  \label{def:physics:phyx-mini-0570:phyxminiproblems-problemphyxmini0570-horizontaldirection}
  \lean{PhyXMiniProblems.ProblemPhyXMini0570.HorizontalDirection}
  Horizontal propagation directions in the one-dimensional diagram
\end{definition}
\begin{proof}
  This is the declared model definition of Horizontal Direction.
\end{proof}
\begin{definition}[Figure Axis Role]
  \label{def:physics:phyx-mini-0570:phyxminiproblems-problemphyxmini0570-figureaxisrole}
  \lean{PhyXMiniProblems.ProblemPhyXMini0570.FigureAxisRole}
  Physical roles of the two plotted axes
\end{definition}
\begin{proof}
  This is the declared model definition of Figure Axis Role.
\end{proof}
\begin{definition}[Figure Text Label]
  \label{def:physics:phyx-mini-0570:phyxminiproblems-problemphyxmini0570-figuretextlabel}
  \lean{PhyXMiniProblems.ProblemPhyXMini0570.FigureTextLabel}
  Literal quantity/text labels visible in image 570.png
\end{definition}
\begin{proof}
  This is the declared model definition of Figure Text Label.
\end{proof}
\begin{definition}[Potential Step Figure]
  \label{def:physics:phyx-mini-0570:phyxminiproblems-problemphyxmini0570-potentialstepfigure}
  \lean{PhyXMiniProblems.ProblemPhyXMini0570.PotentialStepFigure}
  \uses{def:physics:phyx-mini-0570:phyxminiproblems-problemphyxmini0570-horizontaldirection, def:physics:phyx-mini-0570:phyxminiproblems-problemphyxmini0570-figureaxisrole, def:physics:phyx-mini-0570:phyxminiproblems-problemphyxmini0570-figuretextlabel}
  Qualitative evidence read from the primary raster: axes and labels, a zero left plateau, a jump at the origin, the V₀ right plateau, the lower E level, and a particle directed toward positive x
\end{definition}
\begin{proof}
  This is the declared model definition of Potential Step Figure.
\end{proof}
\begin{definition}[Electron Potential Step Setup]
  \label{def:physics:phyx-mini-0570:phyxminiproblems-problemphyxmini0570-electronpotentialstepsetup}
  \lean{PhyXMiniProblems.ProblemPhyXMini0570.ElectronPotentialStepSetup}
  \uses{def:physics:phyx-mini-0570:phyxminiproblems-problemphyxmini0570-lengthquantity, def:physics:phyx-mini-0570:phyxminiproblems-problemphyxmini0570-signedlengthquantity, def:physics:phyx-mini-0570:phyxminiproblems-problemphyxmini0570-massquantity, def:physics:phyx-mini-0570:phyxminiproblems-problemphyxmini0570-inverselengthquantity, def:physics:phyx-mini-0570:phyxminiproblems-problemphyxmini0570-actionquantity, def:physics:phyx-mini-0570:phyxminiproblems-problemphyxmini0570-particlespecies, def:physics:phyx-mini-0570:phyxminiproblems-problemphyxmini0570-potentialregion, def:physics:phyx-mini-0570:phyxminiproblems-problemphyxmini0570-horizontaldirection, def:physics:phyx-mini-0570:phyxminiproblems-problemphyxmini0570-potentialstepfigure}
  Independent quantities in the under-step scattering experiment. In particular, the probability-density penetration distance is stored rather than defined from the source's recorded numerical answer
\end{definition}
\begin{proof}
  This is the declared model definition of Electron Potential Step Setup.
\end{proof}
\begin{definition}[Matches Electron Step Scenario]
  \label{def:physics:phyx-mini-0570:phyxminiproblems-problemphyxmini0570-matcheselectronstepscenario}
  \lean{PhyXMiniProblems.ProblemPhyXMini0570.MatchesElectronStepScenario}
  \uses{def:physics:phyx-mini-0570:phyxminiproblems-problemphyxmini0570-electronpotentialstepsetup}
  The prose describes an electron incident from the left toward positive x
\end{definition}
\begin{proof}
  This is the declared model definition of Matches Electron Step Scenario.
\end{proof}
\begin{definition}[Matches Problem Numerical Readouts]
  \label{def:physics:phyx-mini-0570:phyxminiproblems-problemphyxmini0570-matchesproblemnumericalreadouts}
  \lean{PhyXMiniProblems.ProblemPhyXMini0570.MatchesProblemNumericalReadouts}
  \uses{def:physics:phyx-mini-0570:phyxminiproblems-problemphyxmini0570-energyinelectronvolts, def:physics:phyx-mini-0570:phyxminiproblems-problemphyxmini0570-electronpotentialstepsetup}
  The numerical input data stated in the question
\end{definition}
\begin{proof}
  This is the declared model definition of Matches Problem Numerical Readouts.
\end{proof}
\begin{definition}[Matches Upward Potential Step Profile]
  \label{def:physics:phyx-mini-0570:phyxminiproblems-problemphyxmini0570-matchesupwardpotentialstepprofile}
  \lean{PhyXMiniProblems.ProblemPhyXMini0570.MatchesUpwardPotentialStepProfile}
  \uses{def:physics:phyx-mini-0570:phyxminiproblems-problemphyxmini0570-energyinjoules, def:physics:phyx-mini-0570:phyxminiproblems-problemphyxmini0570-positioninmeters, def:physics:phyx-mini-0570:phyxminiproblems-problemphyxmini0570-electronpotentialstepsetup}
  The sharp step is at the origin, with zero potential on the left and V₀ on the right. The incident total energy equals kinetic plus potential energy
\end{definition}
\begin{proof}
  This is the declared model definition of Matches Upward Potential Step Profile.
\end{proof}
\begin{definition}[Matches Supplied Potential Step Figure]
  \label{def:physics:phyx-mini-0570:phyxminiproblems-problemphyxmini0570-matchessuppliedpotentialstepfigure}
  \lean{PhyXMiniProblems.ProblemPhyXMini0570.MatchesSuppliedPotentialStepFigure}
  \uses{def:physics:phyx-mini-0570:phyxminiproblems-problemphyxmini0570-potentialstepfigure}
  Primary-image evidence from 570.png. This contains no penetration-distance readout and no answer-choice label
\end{definition}
\begin{proof}
  This is the declared model definition of Matches Supplied Potential Step Figure.
\end{proof}
\begin{definition}[Uses Standard Electron Reference Data]
  \label{def:physics:phyx-mini-0570:phyxminiproblems-problemphyxmini0570-usesstandardelectronreferencedata}
  \lean{PhyXMiniProblems.ProblemPhyXMini0570.UsesStandardElectronReferenceData}
  \uses{def:physics:phyx-mini-0570:phyxminiproblems-problemphyxmini0570-massinkilograms, def:physics:phyx-mini-0570:phyxminiproblems-problemphyxmini0570-actioninjouleseconds, def:physics:phyx-mini-0570:phyxminiproblems-problemphyxmini0570-electronpotentialstepsetup}
  Standard electron mass and Physlib reduced-Planck-constant calibrations
\end{definition}
\begin{proof}
  This is the declared model definition of Uses Standard Electron Reference Data.
\end{proof}
\begin{definition}[Has Physical Under Step Parameters]
  \label{def:physics:phyx-mini-0570:phyxminiproblems-problemphyxmini0570-hasphysicalunderstepparameters}
  \lean{PhyXMiniProblems.ProblemPhyXMini0570.HasPhysicalUnderStepParameters}
  \uses{def:physics:phyx-mini-0570:phyxminiproblems-problemphyxmini0570-energyinjoules, def:physics:phyx-mini-0570:phyxminiproblems-problemphyxmini0570-massinkilograms, def:physics:phyx-mini-0570:phyxminiproblems-problemphyxmini0570-actioninjouleseconds, def:physics:phyx-mini-0570:phyxminiproblems-problemphyxmini0570-lengthinmeters, def:physics:phyx-mini-0570:phyxminiproblems-problemphyxmini0570-inverselengthininversemeters, def:physics:phyx-mini-0570:phyxminiproblems-problemphyxmini0570-electronpotentialstepsetup}
  Positivity and E < V₀ select the classically forbidden branch
\end{definition}
\begin{proof}
  This is the declared model definition of Has Physical Under Step Parameters.
\end{proof}
\begin{definition}[Satisfies Evanescent Probability Density Law]
  \label{def:physics:phyx-mini-0570:phyxminiproblems-problemphyxmini0570-satisfiesevanescentprobabilitydensitylaw}
  \lean{PhyXMiniProblems.ProblemPhyXMini0570.SatisfiesEvanescentProbabilityDensityLaw}
  \uses{def:physics:phyx-mini-0570:phyxminiproblems-problemphyxmini0570-energyinjoules, def:physics:phyx-mini-0570:phyxminiproblems-problemphyxmini0570-massinkilograms, def:physics:phyx-mini-0570:phyxminiproblems-problemphyxmini0570-actioninjouleseconds, def:physics:phyx-mini-0570:phyxminiproblems-problemphyxmini0570-inverselengthininversemeters, def:physics:phyx-mini-0570:phyxminiproblems-problemphyxmini0570-electronpotentialstepsetup}
  For a constant region with E < V₀, the amplitude decay constant obeys κ = sqrt (2 m (V₀ - E)) / ℏ, and the probability density relative to its boundary value obeys ρ(x) / ρ(0) = exp (-2 κ x) for x ≥ 0. These are general governing laws, not a numerical penetration-distance answer
\end{definition}
\begin{proof}
  This is the declared model definition of Satisfies Evanescent Probability Density Law.
\end{proof}
\begin{definition}[Is Probability Density Efolding Distance]
  \label{def:physics:phyx-mini-0570:phyxminiproblems-problemphyxmini0570-isprobabilitydensityefoldingdistance}
  \lean{PhyXMiniProblems.ProblemPhyXMini0570.IsProbabilityDensityEfoldingDistance}
  \uses{def:physics:phyx-mini-0570:phyxminiproblems-problemphyxmini0570-lengthinmeters, def:physics:phyx-mini-0570:phyxminiproblems-problemphyxmini0570-electronpotentialstepsetup}
  The requested observable is the positive distance at which the normalized probability density has fallen to exp (-1). This fixes the physical meaning of "penetration distance" without assuming its numerical value
\end{definition}
\begin{proof}
  This is the declared model definition of Is Probability Density Efolding Distance.
\end{proof}
\begin{lemma}[probability Density Penetration Distance closed Form]
  \label{lem:physics:phyx-mini-0570:phyxminiproblems-problemphyxmini0570-probabilitydensitypenetrationdistance-closedform}
  \lean{PhyXMiniProblems.ProblemPhyXMini0570.probabilityDensityPenetrationDistance_closedForm}
  \uses{def:physics:phyx-mini-0570:phyxminiproblems-problemphyxmini0570-energyinjoules, def:physics:phyx-mini-0570:phyxminiproblems-problemphyxmini0570-massinkilograms, def:physics:phyx-mini-0570:phyxminiproblems-problemphyxmini0570-actioninjouleseconds, def:physics:phyx-mini-0570:phyxminiproblems-problemphyxmini0570-lengthinmeters, def:physics:phyx-mini-0570:phyxminiproblems-problemphyxmini0570-electronpotentialstepsetup, def:physics:phyx-mini-0570:phyxminiproblems-problemphyxmini0570-hasphysicalunderstepparameters, def:physics:phyx-mini-0570:phyxminiproblems-problemphyxmini0570-satisfiesevanescentprobabilitydensitylaw, def:physics:phyx-mini-0570:phyxminiproblems-problemphyxmini0570-isprobabilitydensityefoldingdistance}
  The two governing decay laws imply the coherent-SI closed form for the probability-density penetration distance. This is a derived conclusion, not a field of any premise structure
\end{lemma}
\begin{proof}
  The conclusion follows from the governing relations and figure readouts named in the statement of probability Density Penetration Distance closed Form.
\end{proof}
\begin{definition}[Answer Choice]
  \label{def:physics:phyx-mini-0570:phyxminiproblems-problemphyxmini0570-answerchoice}
  \lean{PhyXMiniProblems.ProblemPhyXMini0570.AnswerChoice}
  Labels of the four penetration-distance choices in the problem
\end{definition}
\begin{proof}
  This is the declared model definition of Answer Choice.
\end{proof}
\begin{definition}[displayed Penetration Distance Nanometers]
  \label{def:physics:phyx-mini-0570:phyxminiproblems-problemphyxmini0570-displayedpenetrationdistancenanometers}
  \lean{PhyXMiniProblems.ProblemPhyXMini0570.displayedPenetrationDistanceNanometers}
  \uses{def:physics:phyx-mini-0570:phyxminiproblems-problemphyxmini0570-answerchoice}
  The penetration distance printed beside each label, measured in nanometres
\end{definition}
\begin{proof}
  This is the declared model definition of displayed Penetration Distance Nanometers.
\end{proof}
\begin{definition}[recorded Dataset Answer]
  \label{def:physics:phyx-mini-0570:phyxminiproblems-problemphyxmini0570-recordeddatasetanswer}
  \lean{PhyXMiniProblems.ProblemPhyXMini0570.recordedDatasetAnswer}
  \uses{def:physics:phyx-mini-0570:phyxminiproblems-problemphyxmini0570-answerchoice}
  Dataset metadata recording answer C; it is deliberately not a premise
\end{definition}
\begin{proof}
  This is the declared model definition of recorded Dataset Answer.
\end{proof}
\begin{definition}[Agrees When Rounded To Three Decimal Nanometers]
  \label{def:physics:phyx-mini-0570:phyxminiproblems-problemphyxmini0570-agreeswhenroundedtothreedecimalnanometers}
  \lean{PhyXMiniProblems.ProblemPhyXMini0570.AgreesWhenRoundedToThreeDecimalNanometers}
  \uses{def:physics:phyx-mini-0570:phyxminiproblems-problemphyxmini0570-lengthinnanometers, def:physics:phyx-mini-0570:phyxminiproblems-problemphyxmini0570-electronpotentialstepsetup, def:physics:phyx-mini-0570:phyxminiproblems-problemphyxmini0570-answerchoice, def:physics:phyx-mini-0570:phyxminiproblems-problemphyxmini0570-displayedpenetrationdistancenanometers}
  Agreement with a value displayed to three decimal places in nanometres
\end{definition}
\begin{proof}
  This is the declared model definition of Agrees When Rounded To Three Decimal Nanometers.
\end{proof}
\begin{definition}[Is Unique Nearest Penetration Distance Choice]
  \label{def:physics:phyx-mini-0570:phyxminiproblems-problemphyxmini0570-isuniquenearestpenetrationdistancechoice}
  \lean{PhyXMiniProblems.ProblemPhyXMini0570.IsUniqueNearestPenetrationDistanceChoice}
  \uses{def:physics:phyx-mini-0570:phyxminiproblems-problemphyxmini0570-lengthinnanometers, def:physics:phyx-mini-0570:phyxminiproblems-problemphyxmini0570-electronpotentialstepsetup, def:physics:phyx-mini-0570:phyxminiproblems-problemphyxmini0570-answerchoice, def:physics:phyx-mini-0570:phyxminiproblems-problemphyxmini0570-displayedpenetrationdistancenanometers}
  The selected value is strictly nearer than each other displayed distance
\end{definition}
\begin{proof}
  This is the declared model definition of Is Unique Nearest Penetration Distance Choice.
\end{proof}
% --- Archon declaration topology end ---
% --- Archon physics formalization source end ---
